\chapter{Categories with weak equivalences}
\label{chap:catweq}

ABC cofibration categories can be essentially thought of as 'nicely behaved' categories with weak equivalences. But how much can we say about categories with weak equivalences without bringing in the cofibrations? We will try to find an answer in this chapter.

We will denote ${\bf ho}{\cat M}$ the homotopy category of a category ${\cat M}$ with a class of weak equivalences ${\cat W}$. The homotopy category ${\bf ho}{\cat M}$ is defined by a universal property, but admits a description in terms of generators and relations, starting with the objects and maps of ${\cat M}$ and formally adding inverses of the maps in ${\cat W}$.

The total left (resp. right) derived of a functor $u \colon {\cat M}_1 \ra {\cat M}_2$ between two category pairs $({\cat M}_i, {\cat W}_i)$ for $i = 1, 2$ with localization functors $\gamma_i \colon {\cat M}_i \ra {\bf ho}{\cat M}_i$ is defined as the left (resp. right) Kan extension of $\gamma_2 u$ along $\gamma_1$.

We will use left, resp. right {\it approximation functors} (\definitionref{defn:wleftapproximation}) as a tool for an existence theorem for total derived functors (\theoremref{thm:generalexistenceleftderived}), and a rather technical adjointness property of total derived functors (\theoremref{thm:generalexistencetotalderivedadjoint2}).

The left, resp. right approximation functors $f \colon ({\cat M}^{'}, {\cat W}^{'}) \ra ({\cat M}, {\cat W})$ among other things have the property that they induce an equivalence ${\bf ho}t \colon {\bf ho}{\cat M}^{'} \ra {\bf ho}{\cat M}$. This is proved by the Approximation \theoremref{thm:approxthm}. 

\section{Universes and smallness}
If ${\cat U} \subset {\cat U}^{'}$ are two universes, a $({\cat U}^{'}, {\cat U})$-category ${\cat C}$ has by definition a ${\cat U}^{'}$-small set of objects $Ob{\cat C}$ and ${\cat U}$-small Hom-sets. A category ${\cat C}$ is ${\cat U}$-small if it is a $({\cat U}, {\cat U})$-category, and is {\it locally} ${\cat U}$-small if it has ${\cat U}$-small Hom-sets.

For a fixed universe pair ${\cat U} \subset {\cat U}^{'}$, the ${\cat U}^{'}$-small sets are also referred to as {\it classes}, while the ${\cat U}$-small sets are referred to simply as {\it sets}, or {\it small sets}. We will denote ${\cat Cat}$, resp. $CAT$ the category of ${\cat U}$-small categories, resp. $({\cat U}^{'}, {\cat U})$ categories, with functors as 1-cells. Both ${\cat Cat}$ and $CAT$ actually carry 2-category structures, with natural transformations as 2-cells.

\section{The homotopy category}
\label{sec:whomotopycat}

\begin{defn}
\label{defn:hcat}
Suppose we have a $({\cat U}^{'}, {\cat U})$-category ${\cat M}$ with a class of weak equivalences ${\cat W}$. The homotopy category ${\bf ho}{\cat M}$ by definition is a ${\cat U}^{'}$-small category equipped with a localization functor $\gamma_{\cat M} \colon {\cat M}$ $\ra$ ${\bf ho}{\cat M}$, with the properties that 
\begin{enumerate}
\item
$\gamma_{\cat M}$ sends weak equivalences to isomorphisms
\item
\label{defn:hcat:2}
for any other such functor $\gamma^{'} \colon {\cat M}$ $\ra$ ${\cat M}^{'}$ that sends weak equivalences to isomorphisms, there exists a unique functor $\delta \colon {\bf ho}{\cat M}$ $\ra$ ${\cat M}^{'}$ such that $\delta \gamma_{\cat M} = \gamma^{'}$
\end{enumerate}
\end{defn}

From its universal property, the homotopy category is uniquely defined up to an {\it isomorphism} of categories.

The homotopy category always exists, and is constructed in the next theorem.

\begin{thm}[Gabriel-Zisman]
\label{thm:gz}
For a category with a class of weak equivalences $({\cat M}, {\cat W})$, the category 
\begin{enumerate}
\item
With the same objects as ${\cat M}$ 
\item
With maps between $X$ and $Y$ the equivalence classes of zig-zags 
\[
\xymatrix{ X = A_0 \ar@{-}[r]^-{f_1} &  A_1 \ar@{-}[r]^-{f_2} & ... \ar@{-}[r]^-{f_n} & A_n = Y}
\]
where $f_i$ are maps in ${\cat M}$ going either forward or backward, and all the maps going backward are in ${\cat W}$, where
\begin{enumerate}
\item
Two zig-zags are equivalent if they can be obtained one from another by a finite number of the following three operations and their inverses:
\begin{enumerate}
\item
Skipping elements $A \overset{1_{A}}{\ra} A$ or $A \overset{1_{A}}{\la} A$
\item
Replacing $A \overset{f}{\ra} B \overset{g}{\ra} C$ with $A \overset{gf}{\ra} C$ or $A \overset{f}{\la} B \overset{g}{\la} C$ with $A \overset{fg}{\la} C$
\item
Skipping elements $A \overset{f}{\ra} B \overset{f}{\la} A$ or $A \overset{f}{\la} B \overset{f}{\ra} A$
\end{enumerate}
\item
Composition of maps is induced by the concatenation of zig-zags
\end{enumerate}
\end{enumerate}
is a homotopy category with a localization functor that preserves the objects, and sends a map $X \overset{f}{\ra}Y$ to itself viewed as a zig-zag.
\end{thm}

\begin{proof}
See \cite{Gabriel-Zisman}.
\end{proof}

The homotopy category is sometimes also denoted ${\cat M}[{\cat W}^{-1}]$. 

\index{${\cat W}, \overline{\cat W}$}
\index{weak equivalences ${\cat W}$}
\index{weak equivalences ${\cat W}$!saturation $\overline{\cat W}$}
The saturation $\overline{\cat W}$ of ${\cat W}$ by definition is the class of maps of ${\cat M}$ that become isomorphisms in ${\bf ho}{\cat M}$. The saturation $\overline{\cat W}$ is closed under composition, includes the isomorphisms of ${\cat M}$, and ${\cat M}[\overline{W}^{-1}]$ is isomorphic to ${\cat M}[{\cat W}^{-1}]$.

If ${\cat M}$ is a locally ${\cat U}$-small category with a class of weak equivalences ${\cat W}$, then ${\bf ho}{\cat M} = {\cat M}[{\cat W}^{-1}]$ is not necessarily locally ${\cat U}$-small. For example, if ${\cat M}$ is the $({\cat U}^{'}, {\cat U})$-category of sets ${\cat Set}$, then ${\cat W} = {\cat Cof} = {\cat Fib} = {\cat Set}$  gives ${\cat M}$ a structure of ABC model category, but ${\cat Set}[{\cat Set}^{-1}]$ is only locally ${\cat U}^{'}$-small. 

On the other hand if a locally ${\cat U}$-small category ${\cat M}$ satisfies the hypothesis of \corollaryref{cor:smallhomotopycat} (for example if it is a Quillen model category), then ${\bf ho}{\cat M}$ is again locally ${\cat U}$-small.

One benefit of the intrinsic description of the homotopy category in terms of zig-zags of maps is that it shows that the homotopy category remains the same {\it independent of the universe pair} $({\cat U}^{'}, {\cat U})$ that we start with.

\section{Homotopic maps}
\label{sec:generalhomotopicmaps}

\begin{defn}
\label{defn:homotopicmaps}
For a category ${\cat M}$ with a class of weak equivalences ${\cat W}$, say that two maps $f, g \colon A \ra B$ in ${\cat M}$ are {\it homotopic} (write $f \simeq g$) if $\gamma_{\cat M} f = \gamma_{\cat M} g$.
\end{defn}

The homotopy relation $\simeq$ is an equivalence relation on $Hom_{\cat M}(A, B)$, and the quotient set $Hom_{\cat M}(A, B) / _\simeq$ is classically denoted with the bracket notation $[A, B]$. 

Suppose that $f \simeq g \colon A \ra B$. If $u \colon A^{'} \ra A$, then $fu \simeq gu$. If $u \colon B \ra B^{'}$, then $uf \simeq ug$. 

The quotient category ${\cat M}/_\simeq$, also denoted $\pi {\cat M}$, can therefore be constructed with the same objects as ${\cat M}$ and with $Hom_{\pi {\cat M}} (A, B) = [A, B]$. The induced functor $\pi {\cat M} \ra {\bf ho}{\cat M}$ is an inclusion, and the localization functor $\gamma_{\cat M}$ then factors as the composition ${\cat M} \ra \pi {\cat M} \ra {\bf ho}{\cat M}$.

\section{Left and right approximation functors}
\label{sec:approximationfunctors}

For the rest of the text, it is convenient to restrict ourselves to the case when ${\cat W}$ is closed under composition and includes the identity maps of ${\cat M}$. In the language of \definitionref{defn:catweakequiv}, these are the {\it category pairs} $({\cat M}$, ${\cat W})$. 

A category pair {\it functor} $t \colon ({\cat M}^{'}$, ${\cat W}^{'})$ $\ra$ $({\cat M}$, ${\cat W})$ by definition will be a functor $t \colon {\cat M}^{'}$ $\ra$ ${\cat M}$ that preserves the weak equivalences. A category pair functor $t$ as above induces a functor at the level of the homotopy categories, denoted ${\bf ho}t \colon {\bf ho}{\cat M}^{'}$ $\ra$ ${\bf ho}{\cat M}$.

Our next goal is to find sufficient conditions for a category pair functor to induce an {\it equivalence} at the level of homotopy categories.

\index{approximation!left, right}
\begin{defn}
\label{defn:wleftapproximation}
A category pair functor $t \colon ({\cat M}^{'}$, ${\cat W}^{'})$ $\ra$ $({\cat M}$, ${\cat W})$ is a {\it left approximation} if
\begin{description}
\item[LAP1]
For any object $A \eps {\cat M}$ there exists an object $A^{'} \eps {\cat M}^{'}$ and a weak equivalence map $tA^{'} \overset{\sim}{\ra} A$
\item[LAP2]
For any maps $f \eps {\cat M}$ and $p \eps {\cat W}$ as below
\[
\xymatrix@R=15pt@C=15pt { 
&&&& B \\
&&&& \\
tA^{'} \ar[rrrruu]^{f} \ar@{-->}[rr]_{ti^{'}} && tD^{'} \ar@{-->}[rruu]^{\sigma} && tB^{'} \ar@{-->}[ll]^{tj^{'}}_\sim \ar[uu]_{p}^\sim
}
\]
\begin{enumerate}
\item
There exist maps $A^{'} \overset{i^{'}}{\lra} D^{'} \underset{\sim}{\overset{j^{'}}{\lla}} B^{'}$ with $j^{'} \eps {\cat W}^{'}$ and a weak equivalence $\sigma \eps {\cat W}$ making the diagram commutative.
\item
For any other such $i^{''}, j^{''}, \sigma^{''}$ the equality $j^{'-1}i^{'} = j^{''-1}i^{''}$ holds in ${\bf ho}{\cat M}^{'}$.
\end{enumerate}
%\[
%tA^{'} \overset{f}{\lra} B \underset{\sim}{\overset{s}{\lla}} tC^{'}
%\]
\end{description}

Dually, the functor $t$ is a {\it right approximation} if
\begin{description}
\item[RAP1]
For any object $A \eps {\cat M}$ there exists an object $A^{'} \eps {\cat M}^{'}$ and a weak equivalence map $A \overset{\sim}{\ra} tA^{'}$
\item[RAP2]
For any maps $f \eps {\cat M}$ and $i \eps {\cat W}$ as below
\[
\xymatrix@R=15pt@C=15pt { 
A \ar[dd]_{i}^\sim \ar@{-->}[rrdd]^{\tau} \ar[rrrrdd]^{f}  &&&& \\
&&&& \\
tA^{'}  && tS^{'} \ar@{-->}[ll]^{tq^{'}}_\sim \ar@{-->}[rr]_{tp^{'}} && tB^{'}  
}
\]
\begin{enumerate}
\item
There exist maps $A^{'} \underset{\sim}{\overset{q^{'}}{\lla}} S^{'} \overset{p^{'}}{\lra} B^{'}$ with $q^{'} \eps {\cat W}^{'}$ and a weak equivalence $\tau \eps {\cat W}$ making the diagram commutative.
\item
For any other such $p^{''}, q^{''}, \tau^{''}$ the equality $p^{'}q^{'-1} = p^{''}q^{''-1}$ holds in ${\bf ho}{\cat M}^{'}$.
\end{enumerate}
\end{description}
\end{defn}

%\begin{defn}
%\label{defn:absoluteapprox}
%A left (resp. right) approximation $t \colon ({\cat M}^{'}, {\cat W}^{'}) \ra ({\cat M}, {\cat W})$ is {\it absolute} if for any small category ${\cat D}$, the induced functor of diagram categories $(({\cat M}^{'})^{\cat D}, ({\cat W}^{'})^{\cat D}) \ra ({\cat M}^{\cat D}, {\cat W}^{\cat D})$ is a left (resp. right) approximation.
%\end{defn}

\begin{lem}
\label{lem:sufficientcondequiv}
Suppose that $t \colon ({\cat M}^{'}, {\cat W}^{'}) \ra ({\cat M}, {\cat W})$ is a category pair functor, and $f \colon A^{'} \ra C^{'}$ is a map in ${\cat M}^{'}$. Suppose that either
\begin{enumerate}
\item
$t$ is a left approximation, and there exist $p, q \eps {\cat W}$ making commutative the diagram
\[
\xymatrix{
  & B & \\
  tA^{'} \ar[ru]^p_\sim \ar[rr]_{tf^{'}} && tC^{'} \ar[lu]^\sim_q
}
\]
\item
$t$ is a right approximation, and there exist $i, j \eps {\cat W}$ making commutative the diagram
\[
\xymatrix{
  & B \ar[ld]_i^\sim \ar[rd]^j_\sim & \\
  tA^{'} \ar[rr]_{tf^{'}} && tC^{'} 
}
\]
\end{enumerate}
Then the image of $f^{'}$ in ${\bf ho}{\cat M}^{'}$ is an isomorphism, i.e. $f \eps \overline{\cat W}^{'}$.
\end{lem}

\begin{proof}
We only prove (1). Using LAP2 (1), we can construct $A^{'} \underset{\sim}{\overset{i^{'}}{\lra}} B^{'} \overset{j^{'}}{\lla} C^{'}$ with $i^{'} \eps {\cat W}^{'}$ and a weak equivalence $\sigma \eps {\cat W}$ with $\sigma \comp ti^{'} = p$ and $\sigma \comp tj^{'} = q$. 
\[
\xymatrix {
& B & \\ 
tA^{'} \ar[rr]_(.3){tf^{'}}|{\;\;\;} \ar[ru]^p \ar@{-->}[rd]_{ti^{'}} && tC^{'} \ar@{-->}[ld]^{tj^{'}} \ar[lu]_q \\
& tB^{'} \ar@{-->}[uu]^(.7){\sigma} & 
} 
\]
From LAP2 (2), $j^{'}f^{'}$ is invertible in ${\bf ho}{\cat M}^{'}$. Similarly we can construct $j^{''} \eps {\cat M}^{'}$ such that $j^{''}j^{'}$ is invertible in ${\bf ho}{\cat M}^{'}$. We deduce that $j^{'}$ and therefore $f^{'}$ are invertible in ${\bf ho}{\cat M}^{'}$.
\end{proof}

\index{cleavage!left (right) cleavage}
\begin{defn}
\label{defn:wleftcleavage}
Let $t \colon {\cat M}^{'} \ra {\cat M}$ be a left approximation of category pairs $({\cat M}^{'}$, ${\cat W}^{'})$ and $({\cat M}$, ${\cat W})$. A {\it left cleavage} ${\cat C} = (C, D, p, i, j, \sigma)$ along $t$ consists of the following data:
\begin{enumerate}
\item
For any object $A$ of ${\cat M}$, an object $C(A)$ of ${\cat M}^{'}$ and a weak equivalence $p_A \colon tC(A) \ra A$
\item
For any map $f \colon A \ra B$, maps $C(A) \overset{i(f)}{\lra} D(f) \underset{\sim}{\overset{j(f)}{\lla}} C(B)$ and a commutative diagram
\[
\xymatrix{
  A \ar[rr]^{f} && B \\
  tC(A) \ar[u]^{p_A}_\sim \ar[r]_{ti(f)} & tD(f) \ar[ur]^{\sigma(f)} & tC(B) \ar[l]^{tj(f)}_\sim \ar[u]^\sim_{p_B}
}
\]
with $j(f)$, $\sigma(f)$ weak equivalences.
\end{enumerate}

A {\it left cleavage map} ${\cat C}_1 \overset{\cat F}{\ra} {\cat C}_2$ from ${\cat C}_1$ $=$ ($C_1$, $D_1$, $p_1$, $i_1$, $j_1$, $\sigma_1$) to ${\cat C}_2$ $=$ ($C_2$, $D_2$, $p_2$, $i_2$, $j_2$, $\sigma_2$) consists of a family of maps $c(A) \colon C_1(A) \ra C_2(A)$ and $d(f) \colon D_1(f) \ra D_2(f)$, satisfying $p_{1A} = p_{2A} \comp tc(A)$, $d(f)i_1(f) = i_2(f)c(A)$ and $d(f)j_1(f) = j_2(f)c(B)$. 

Notice that by \lemmaref{lem:sufficientcondequiv}, we must have $c(A), d(f) \eps \overline{\cat W}^{'}$.

Composition of left cleavage maps is given by componentwise composition of $c(A), d(f)$. Left cleavages and left cleavage maps along $t$ form therefore a category.
\end{defn}

The dual definition for right approximations is

\begin{defn}
\label{defn:wrightcleavage}
Let $t \colon {\cat M}^{'} \ra {\cat M}$ be a right approximation of category pairs $({\cat M}^{'}$, ${\cat W}^{'})$ and $({\cat M}$, ${\cat W})$. A {\it right cleavage} ${\cat R} = (R, S, i, p^{'}, q^{'}, \tau)$ along $t$ consists of the following data:
\begin{enumerate}
\item
For any object $A$ of ${\cat M}$, an object $R(A)$ of ${\cat M}^{'}$ and a weak equivalence $i_A \colon A \ra tR(A)$
\item
For any map $f \colon A \ra B$, maps $R(A) \underset{\sim}{\overset{q(f)}{\lla}} S(f) \overset{p(f)}{\lra} R(B)$ and a commutative diagram
\[
\xymatrix{
  A \ar[rr]^{f} \ar[d]_{i_A}^\sim \ar[dr]^{\tau(f)} && B \ar[d]_\sim^{i_B}   \\
  tR(A)  & tS(f)  \ar[l]^{tq(f)}_\sim \ar[r]_{tp(f)} & tR(B) 
}
\]
with $q(f)$, $\tau(f)$ weak equivalences.
\end{enumerate}

A {\it right cleavage map} ${\cat R}_1 \overset{\cat F}{\ra} {\cat R}_2$ from ${\cat R}_1$ $=$ ($R_1$, $S_1$, $i_1$, $p^{'}_1$, $q^{'}_1$, $\tau_1$) to ${\cat R}_2$ = ($R_2$, $S_2$, $i_2$, $p^{'}_2$, $q^{'}_2$, $\tau_2$) consists of a family of maps $r(A) \colon R_1(A) \ra R_2(A)$ and $s(f) \colon S_1(f) \ra S_2(f)$, satisfying $r(A) \comp i_{1A} = i_{2A}$, $r(B)p_1(f) = p_2(f)s(f)$ and $r(A)q_1(f) = q_2(f)s(f)$. 
\end{defn}


\begin{lem}
\label{lem:cleavagemap}
Suppose that $t \colon ({\cat M}^{'}, {\cat W}^{'}) \ra ({\cat M}, {\cat W})$ is a category pair functor.
\mbox{}
\begin{enumerate}
\item
If $t$ is a left approximation and ${\cat C}_1, {\cat C}_2$ are left cleavages, there exist a left cleavage ${\cat C}_3$ and left cleavage maps ${\cat C}_1 \ra {\cat C}_3 \la {\cat C}_2$.
\item
If $t$ is a right approximation and ${\cat R}_1, {\cat R}_2$ are right cleavages, there exist a right cleavage ${\cat R}_3$ and right cleavage maps ${\cat R}_1 \la {\cat R}_3 \ra {\cat R}_2$.
\end{enumerate}
\end{lem}

\begin{proof}
Apply repeatedly the axioms LAP2 (1), resp. RAP2 (1).
\end{proof}

\index{cleavage!normalized}
\index{cleavage!regular object}
\index{cleavage!regular}
\begin{defn}
\label{defn:wregularcleavageobj}
Suppose that $({\cat M}^{'}$, ${\cat W}^{'})$ and $({\cat M}$, ${\cat W})$ are two category pairs.
\begin{enumerate}
\item
Let $t \colon {\cat M}^{'} \ra {\cat M}$ be a left approximation functor, and ${\cat C} = (C, D, p, i, j, \sigma)$ be a left cleavage. 
\begin{enumerate}
\item
The left cleavage ${\cat C}$ is {\it normalized} if for any object $A$ of ${\cat M}$ we have $D(1_A) = C(A)$, $i(1_A) = j(1_A) = 1_{C(A)}$ and $\sigma(1_A) = p_A$.
\item
An object $A^{'} \eps {\cat M}^{'}$ is {\it regular} with respect to ${\cat C}$ (referred to simply as {\it regular} if no confusion is possible) if $C(tA^{'}) = A^{'}$ and $p_{tA^{'}} = 1_{tA^{'}}$.
\item
The left cleavage ${\cat C}$ is {\it regular} if it is normalized, and any object $A^{'} \eps {\cat M}^{'}$ admits a regular object $A^{''} \eps {\cat M}^{'}$ isomorphic to $A^{'}$ in ${\bf ho}{\cat M}^{'}$.
\end{enumerate}
\item
Let $t \colon {\cat M}^{'} \ra {\cat M}$ be a right approximation functor, and ${\cat R} = (R, S, i, p^{'}, q^{'}, \tau)$ be a right cleavage along $t$. 
\begin{enumerate}
\item
The right cleavage ${\cat R}$ is {\it normalized} if for any object $A$ of ${\cat M}$ we have $S(1_A) = R(A)$, $p(1_A) = q(1_A) = 1_{R(A)}$ and $\tau(1_A) = i_A$.
\item
An object $A^{'} \eps {\cat M}^{'}$ is {\it regular} with respect to a right cleavage along $t$ if $R(tA^{'}) = A^{'}$ and $i_{tA^{'}} = 1_{tA^{'}}$.
\item
The right cleavage ${\cat R}$ is {\it regular} if it is normalized, and any object $A^{'} \eps {\cat M}^{'}$ admits a regular object $A^{''} \eps {\cat M}^{'}$ isomorphic to $A^{'}$ in ${\bf ho}{\cat M}^{'}$.
\end{enumerate}
\end{enumerate}
\end{defn}

\begin{lem}
\label{lem:admitframing}
\mbox{}
\begin{enumerate}
\item
Any left approximation admits a regular left cleavage.
\item
Any right approximation admits a regular right cleavage. 
\end{enumerate}
\end{lem}

\begin{proof}
We only prove (1), and perform the construction in a number of steps.

(i) For any object $A \eps {\cat M}$ of the form $A = tA^{'}$, pick exactly one such $A^{'}$ and define $C(tA^{'}) = A^{'}$, $p_{tA^{'}} = 1_{tA^{'}}$.

(ii). For any $A \eps {\cat M}$ for which this data is not defined, use LAP1 to define the object $C(A)$ and the weak equivalence $p_A$.

(iii). For any $A \eps {\cat M}$, define $D(1_A) = C(A)$, $i(1_A) = j(1_A) = 1_{C(A)}$ and $\sigma(1_A) = p_A$.

(iv). For any map $f \colon A \ra B$ in ${\cat M}$ for which this data is not defined, use LAP2 (1) to construct the desired objects and maps $D(f), i(f) \eps {\cat M^{'}}$, $j(f) \eps {\cat W}^{'}$ and $\sigma(f) \eps {\cat W}$.

This constructs a left cleavage, which is normalised because of step (iii). 

Let us check condition (1) of \definitionref{defn:wregularcleavageobj}. For $A^{'} \eps {\cat M}^{'}$, we have $tC(tA^{'}) = tA^{'}$. The object $C(tA^{'})$ is regular, and by \lemmaref{lem:sufficientcondequiv} it is isomorphic to $A^{'}$ in ${\bf ho}{\cat M}$. In conclusion, we have constructed a regular left cleavage.
\end{proof}

If a left (resp. right) approximation $t \colon {\cat M}^{'} \ra {\cat M}$ is {\it injective on objects}, we could actually construct a normalized left (resp. right) cleavage where all objects of ${\cat M}^{'}$ were regular. 

\section{The approximation theorem}

Recall that a functor $u \colon {\cat M}_1 \ra {\cat M}_2$ is called
\index{functor!full}
\index{functor!faithful}
\index{functor!essentially surjective}
\begin{enumerate}
\item
{\it essentially surjective} if any object of ${\cat M}_2$ is isomorphic to an object in the image of $u$
\item
{\it full} if any map in $Hom_{{\cat M}_2}(uA, uB)$ is in the image of $u$, for all objects $A, B$ of ${\cat M}_1$
\item
{\it faithful} if $u$ is injective on $Hom_{{\cat M}_1}(A, B)$ for all objects $A, B$ of ${\cat M}_1$
\end{enumerate}

The functor $u$ is an equivalence of categories if and only if it is essentially surjective, full and faithful.


\begin{thm}[The approximation theorem]
\label{thm:approxthm}
\mbox{} \\
A left (or right) approximation functor $t \colon ({\cat M}^{'}$, ${\cat W}^{'})$ $\ra$ $({\cat M}$, ${\cat W})$ induces an equivalence of homotopy categories ${\bf ho} t \colon{\bf ho}{\cat M}^{'} \ra {\bf ho}{\cat M}$.
\end{thm}

\begin{proof}
Suppose that $t$ is a left approximation. For a left cleavage ${\cat C} = (C, D, p, i, j, \sigma)$ along $t$, construct a functor ${\bf s}_{\cat C} \colon {\bf ho}{\cat M}^{'} \ra {\bf ho}{\cat M}$, as follows:
\begin{enumerate}
\item
On objects $A \eps {\cat M}$, define ${\bf s}_{\cat C}A  = C(A)$
\item
For a zig-zag ${\cat Z}$ in ${\cat M}$
\begin{equation}
\label{eqn:approxarehtpyequiv}
\xymatrix{ A = A_0 \ar@{-}[r]^-{f_1} &  A_1 \ar@{-}[r]^-{f_2} & ... \ar@{-}[r]^-{f_n} & A_n = B}
\end{equation}
where the maps $f_k$ go either forward or backward, and all the maps going backward are in ${\cat W}$, denote $A^{'}_{k} = C(A_k)$ and $p_k = p_{A_k}$.
\begin{enumerate}
\item
If an $f_k$ goes forward, then construct the commutative diagram
\begin{equation}
\label{eqn:approxarehtpyequiv1}
\xymatrix{
  A_{k-1} \ar[rr]^{f_k} && A_k \\
  tA^{'}_{k-1} \ar[u]^{p_{k-1}}_\sim \ar[r]_{ti^{'}_k} & tD^{'}_k \ar[ur]^{\sigma_k} & tA^{'}_k \ar[l]^{tj^{'}_k}_\sim \ar[u]^\sim_{p_k}
}
\end{equation}
\item
If an $f_k$ goes backward, then construct the commutative diagram
\begin{equation}
\label{eqn:approxarehtpyequiv2}
\xymatrix{
  A_{k-1}  && A_k \ar[ll]^\sim_{f_k} \\
  tA^{'}_{k-1} \ar[u]^{p_{A_{k-1}}}_\sim \ar[r]_{ti^{'}_k} & tD^{'}_k \ar[ul]_{\sigma_k} & tA^{'}_k \ar[l]^{tj^{'}_k}_\sim \ar[u]^\sim_{p_k}
}
\end{equation}
\end{enumerate}
where in both cases we denoted $D^{'}_k = D(f_k)$, $i^{'}_k = i(f_k)$, $j^{'}_k = j(f_k)$ and $\sigma_k = \sigma(f_k)$.
\end{enumerate}

The maps $i^{'}_k, j^{'}_k$ collect together to a zig-zag denoted ${\cat Z}^{'}$ in ${\cat M}^{'}$, from $A^{'}_0$ to $A^{'}_n$, with all backward going maps being weak equivalences. We define ${\bf s}_{\cat C}$(image of ${\cat Z}$ in ${\bf ho}{\cat M}$) = image of ${\cat Z}^{'}$ in ${\bf ho}{\cat M}^{'}$.

Let's show that the definition of ${\bf s}_{\cat C}$ on maps does not depend on the choices involved. Let ${\cat Z}_1$, ${\cat Z}_2$ by two zig-zags of the form \equationref{eqn:approxarehtpyequiv}, and denote ${\cat Z}^{'}_1$, ${\cat Z}^{'}_2$ the associated zig-zags constructed in ${\cat M}^{'}$.

(i). If ${\cat Z}_2$ is obtained from ${\cat Z}_1$ by inserting an element $A_k \overset{1_{A_k}}{\lra} A_k$, then ${\cat Z}^{'}_2$ is obtained from ${\cat Z}^{'}_1$ by inserting an element 
\[
A_k^{'} \overset{i^{'}}{\lra} D^{'} \underset{\sim}{\overset{j^{'}}{\lla}} A^{'}_k
\]
with the property that there exists a commutative diagram
\[
\xymatrix{
  A_k \ar[rr]^{1_{A_k}} && A_k \\
  tA^{'}_k \ar[u]^{p_k}_\sim \ar[r]_{ti^{'}} & tD^{'} \ar[ur]^{\sigma} & tA^{'}_k \ar[l]^{tj^{'}}_\sim \ar[u]^\sim_{p_k}
}
\]
with $\sigma \eps {\cat W}$. Using LAP2 (2) we see that ${\cat Z}^{'}_1, {\cat Z}^{'}_2$ define the same element in ${\bf ho}{\cat M}^{'}$.

(ii). If ${\cat Z}_2$ is obtained from ${\cat Z}_1$ by inserting an element $A_k \overset{1_{A_k}}{\lla} A_k$, then similar to (i) we can show that ${\cat Z}^{'}_1, {\cat Z}^{'}_2$ define the same element in ${\bf ho}{\cat M}^{'}$.

(iii).  If ${\cat Z}_2$ is obtained from ${\cat Z}_1$ by replacing an element $A_{k-1} \overset{f_k}{\lra} A_{k} \overset{f_{k+1}}{\lra} A_{k+1}$ with $A_{k-1} \overset{f_{k+1}f_k}{\lra} A_{k+2}$ , then ${\cat Z}^{'}_2$ is obtained from ${\cat Z}^{'}_1$ by replacing the element 
\[
A_{k-1}^{'} \overset{i^{'}_{k}}{\lra} D^{'}_{k} \underset{\sim}{\overset{j^{'}_{k}}{\lla}} A^{'}_{k} \overset{i^{'}_{k+1}}{\lra} D^{'}_{k+1} \underset{\sim}{\overset{j^{'}_{k+1}}{\lla}} A^{'}_{k+1}
\]
with an element
\[
A_{k-1}^{'} \overset{i^{'}}{\lra} D^{'} \underset{\sim}{\overset{j^{'}}{\lla}} A^{'}_{k+1}
\]
with the property that there exists a commutative diagram
\[
\xymatrix{
  A_{k-1} \ar[rr]^{f_{k+1}f_k} && A_{k+1} \\
  tA^{'}_{k-1} \ar[u]^{p_{k-1}}_\sim \ar[r]_{ti^{'}} & tD^{'} \ar[ur]^{\sigma} & tA^{'}_{k+1} \ar[l]^{tj^{'}}_\sim \ar[u]^\sim_{p_{k+2}}
}
\]
with $\sigma \eps {\cat W}$. Using LAP2 (1) construct $D^{''}, i^{''} \eps {\cat M}^{'}$, $j^{''} \eps {\cat W}^{''}$ and $\sigma^{''} \eps {\cat W}$ making commutative the diagram
\[
\xymatrix {
& A_{k+1} & \\ 
\; tD^{'}_{k} \; \ar[ru]^{f_{k+1}\sigma_{k}} \ar@{-->}[rd]_{ti^{''}} && tD^{'}_{k+1} \ar@{-->}[ld]^{tj^{''}}_\sim \ar[lu]_{\sigma_{k+1}}^\sim \\
& tD^{''} \ar@{-->}[uu]^{\sigma^{''}} & 
} 
\]

In the diagram
\[
\xymatrix {
  && tA^{'}_k \ar[ld]_{tj^{'}_k}^\sim \ar[rd]^{ti^{'}_{k+1}} &&& \\
  tA^{'}_{k-1} \ar[r]^{ti^{'}_{k}} \ar[rrdd]_(.3){ti^{'}} & tD^{'}_{k} \ar[rd]_{ti^{''}} && tD^{'}_{k+1} \ar[ld]^{tj^{''}}_\sim & tA^{'}_{k+1} \ar[l]_{tj^{'}_{k+1}}^\sim \ar[lldd]^(.3){tj^{'}}_(.3)\sim & \\
  && tD^{''} \ar@{-->}[rrrd]^(.7){\sigma^{''}} &&&  \\
  && tD^{'} \ar@{-->}[rrr]_(.7){\sigma} &&& A_{k+1}
} 
\]
we have $\sigma^{''} \comp t(i^{''}j^{'}_k) = \sigma^{''} \comp t(j^{''}i^{'}_{k+1})$, $\sigma^{''} \comp t(i^{''}i^{'}_k) = \sigma \comp ti^{'}$ and $\sigma^{''} \comp t(j^{''}j^{'}_{k+1}) = \sigma \comp tj^{'}$. Using LAP2 (2) we see that ${\cat Z}^{'}_1, {\cat Z}^{'}_2$ define the same element in ${\bf ho}{\cat M}^{'}$.

(iv).  If ${\cat Z}_2$ is obtained from ${\cat Z}_1$ by replacing weak equivalences $A_{k-1} \underset{\sim}{\overset{f_k}{\lla}} A_{k} \underset{\sim}{\overset{f_{k+1}}{\lla}} A_{k+1}$ with $A_{k-1} \underset{\sim}{\overset{f_{k+1}f_k}{\lla}} A_{k+1}$ , then similar to (iii) one shows that ${\cat Z}^{'}_1, {\cat Z}^{'}_2$ define the same element in ${\bf ho}{\cat M}^{'}$.

(v). If ${\cat Z}_2$ is obtained from ${\cat Z}_1$ by inserting weak equivalences $A_{k} \underset{\sim}{\overset{f}{\lra}} A \underset{\sim}{\overset{f}{\lla}} A_{k}$, then ${\cat Z}^{'}_2$ is obtained from ${\cat Z}^{'}_1$ by inserting an element
\[
A_{k}^{'} \overset{i^{'}}{\lra} D^{'} \underset{\sim}{\overset{j^{'}}{\lla}} A^{'} \underset{\sim}{\overset{j^{'}}{\lra}} D^{'} \overset{i^{'}}{\lla}  A^{'}_k
\]
with the property that there exists a commutative diagram
\[
\xymatrix{
  A_{k}  \ar[rr]^f_\sim && A && A_{k} \ar[ll]_{f}^\sim \\
  tA^{'}_{k} \ar[u]^{p_{k}}_\sim \ar[r]_{ti^{'}} & tD^{'} \ar[ur]^{\sigma} & tA^{'} \ar[l]^{tj^{'}}_\sim \ar[r]_{tj^{'}}^\sim  \ar[u]^\sim_{p} & tD^{'} \ar[lu]_\sigma & tA^{'}_{k} \ar[u]_{p_{k}}^\sim \ar[l]^{ti^{'}} 
}
\]
with $\sigma \eps {\cat W}$. Using \lemmaref{lem:sufficientcondequiv}, we see that $i^{'} \eps \overline{\cat W}^{'}$, and so ${\cat Z}^{'}_1, {\cat Z}^{'}_2$ define the same element in ${\bf ho}{\cat M}^{'}$.

(vi). If ${\cat Z}_2$ is obtained from ${\cat Z}_1$ by replacing weak equivalences $A_k$ with $A_{k} \underset{\sim}{\overset{f}{\lla}} A \underset{\sim}{\overset{f}{\lra}} A_{k}$, similar to (v) one can show that ${\cat Z}^{'}_1, {\cat Z}^{'}_2$ define the same element in ${\bf ho}{\cat M}^{'}$.

From (i)-(vi) we conclude that ${\bf s}_{\cat C}$ is well defined on the maps of ${\bf ho}{\cat M}$.

(vii). If the end of ${\cat Z}_1$ and the beginning of ${\cat Z}_2$ coincide, then so do the end of ${\cat Z}^{'}_1$ and the beginning of ${\cat Z}^{'}_2$. From this observation, we see that ${\bf s}_{\cat C}$ preserves composition of maps.

(viii). A proof similar to (i) shows that ${\bf s}_{\cat C}$ preserves unit morphisms.

Using (vii) and (viii), we see that ${\bf s}_{\cat C} \colon {\bf ho}{\cat M} \ra {\bf ho}{\cat M}^{'}$ is a functor.

The functor ${\bf s}_{\cat C}$ is essentially surjective on objects - this can be verified using LAP1, LAP2 and \lemmaref{lem:sufficientcondequiv}.

From the commutativity of \equationref{eqn:approxarehtpyequiv1} and \equationref{eqn:approxarehtpyequiv2}, we get a natural isomorphism $({\bf ho} t) \comp {\bf s}_{\cat C} \cong 1_{{\bf ho}{\cat M}}$. This shows that ${\bf s}_{\cat C}$ is faithful. 

The functor ${\bf s}_{\cat C}$ is also full - to see that, pick a zig-zag ${\cat Z}^{''}$ in ${\cat M}^{'}$
\[
\xymatrix{ B^{'}_0 \ar@{-}[r]^-{g^{'}_1} &  B^{'}_1 \ar@{-}[r]^-{g^{'}_2} & ... \ar@{-}[r]^-{g^{'}n} & B^{'}_n}
\]
and denote $A_k = tB^{'}_k$, $f_k = tg^{'}_k$. The equivalence class of the zig-zag ${\cat Z}$ of \equationref{eqn:approxarehtpyequiv} is mapped by ${\bf s}_{\cat C}$ to the equivalence class of the zig-zag ${\cat Z}^{'}$ described earlier. For each $k$, using LAP2 (1) we can construct objects $A^{''}_k \eps {\cat M}^{'}$, maps $p^{''}_k \colon A^{''}_k \ra  B^{'}_k$, $q^{''}_k \colon A^{''}_k \ra  A^{'}_k$ and $p^{'}_k \colon tA^{''}_k \ra A_k$  such that $q^{''}_k \eps {\cat W}^{'}$, $p^{'}_k \eps {\cat W}$, $p^{'}_k \comp tp^{''}_k = 1_{A_k}$, $p^{'}_k \comp tq^{''}_k = p_k$. The maps $p^{''}_k$ are in $\overline{\cat W}^{'}$ by \lemmaref{lem:sufficientcondequiv}. 

Using LAP2 (1) twice in a row we can then construct objects $D^{''}_k \eps {\cat M}^{'}$ and maps $i^{''}_k \colon A^{''}_{k-1} \ra D^{''}_k$, $j^{''} \colon A^{''}_k \ra D^{''}_k$, $\nu^{'}_k \colon D^{'}_k \ra D^{''}_k$ and $\tau_k \colon tD^{''}_k \ra (\mathrm{\; target \; of \; }f_k)$ such that $j^{''}_k \eps {\cat W}^{'}$, $\nu^{'}_k \eps \overline{\cat W}^{'}$ with $\nu^{'}_k i^{'}_k q^{''}_k = i^{''}_k$, $\nu^{'}_k j^{'}_k q^{''}_k = j^{''}_k$ and $\tau_k \comp t \nu^{'}_k = \sigma_k$. 

The maps $p^{''}_k, q^{''}_k, i^{''}_k, j^{''}_k $ yield a zig-zag ${\cat Z}^{'''}$. Using LAP2 (2), we see that ${\cat Z}^{'}$ and ${\cat Z}^{'''}$, resp. ${\cat Z}^{''}$ and ${\cat Z}^{'''}$ define isomorphic maps in ${\bf ho}{\cat M}^{'}$. This concludes the proof that ${\bf s}_{\cat C}$ is full.

In conclusion, ${\bf s}_{\cat C}$ is an equivalence of categories, and therefore so is its quasi-inverse ${\bf ho}t$.
\end{proof}

\begin{rem}
Since ${\bf ho} t$ is an equivalence, for any other left cleavage ${\cat C}^{'}$ we have a canonical isomorphism ${\bf s}_{\cat C} \cong {\bf s}_{{\cat C}^{'}}$. By \lemmaref{lem:cleavagemap}, there exist a left cleavage ${\cat C}^{''}$ and a zig-zag of left cleavage maps ${\cat C} \ra {\cat C}^{''} \la {\cat C}^{'}$. It is straightforward to see that any such zig-zag induces the canonical isomorphisms ${\bf s}_{\cat C} \cong {\bf s}_{{\cat C}^{''}} \cong {\bf s}_{{\cat C}^{'}}$.
\end{rem}

\begin{cor}
\label{cor:equivleftapprox}
Suppose that $t \colon ({\cat M}^{'}$, ${\cat W}^{'})$ $\ra$ $({\cat M}$, ${\cat W})$ is a category pair functor.
\begin{enumerate}
\item
The following statements are equivalent:
\begin{enumerate}
\item
$t$ is a left approximation
\item
$t$ satisfies LAP1, LAP2 (1) and induces an equivalence of categories ${\bf ho}t$
\end{enumerate}
\item
The following statements are equivalent:
\begin{enumerate}
\item
$t$ is a right approximation
\item
$t$ satisfies RAP1, RAP2 (1) and induces an equivalence of categories ${\bf ho}t$
\end{enumerate}
\end{enumerate}
\end{cor}

\begin{proof}
(b) $\Ra$ (a) is immediate, and (a) $\Ra$ (b) follows from the Approximation \theoremref{thm:approxthm}.
\end{proof}

\begin{cor}
\label{cor:otherleftapprox}
Left (resp. right) approximation functors are closed under composition.
\end{cor}

\begin{proof}
Category pair functors satisfying LAP1 and LAP2 (1), resp. RAP1 and RAP2 (1) are closed under composition. The corollary now is a consequence of \corollaryref{cor:equivleftapprox}.
\end{proof}

\section{Total derived functors}
\label{sec:totalleftrightderived}

A functor $u \colon {\cat M}_1 \ra {\cat M}_2$ between two category pairs $({\cat M}_1, {\cat W}_1), ({\cat M}_2, {\cat W}_2)$ descends to a functor ${\bf ho} u \colon {\bf ho} {\cat M}_1 \ra {\bf ho} {\cat M}_2$ if and only if $u({\cat W}_1) \subset \overline{\cat W}_2$, where $\overline{\cat W}_2$ denotes the saturation of ${\cat W}_2$ in ${\cat M}_2$.

In the general case however ${\bf ho} u$ does not exist, and the best we can hope for is the existence of a left (or a right) Kan extension of $\gamma_{{\cat M}_2}u$ along $\gamma_{{\cat M}_1}$, also called the {\it total left} (resp. right) {\it derived functors} of $u$.

\index{functor!total derived}
\begin{defn}
\label{defn::def:totalleftkanextensions}
Suppose that $({\cat M}_1, {\cat W}_1), ({\cat M}_2, {\cat W}_2)$ are two categories with weak equivalences, with localization functors denoted $\gamma_{{\cat M}_1}$ and respectively $\gamma_{{\cat M}_1}$, and suppose that $u \colon {\cat M}_1 \ra {\cat M}_2$ is a functor.
\begin{enumerate}
\item
The {\it total left derived} functor of $u$, denoted $({\bf L}u, \eps)$ is the left Kan extension $({\bf L}_{\gamma_{{\cat M}_1}} (\gamma_{{\cat M}_2} u), \eps)$ of $\gamma_{{\cat M}_2} u$ along $\gamma_{{\cat M}_1}$
\begin{center}
$\xymatrix {
    {\cat M}_1 \ar[r]^u \ar[d]_{\gamma_{_{{\cat M}_1}}} & 
    {\cat M}_2 \ar[d]^{\gamma_{_{{\cat M}_2}}} \\
    {\bf ho}{\cat M}_1 \ar[r]_{{\bf L}u}  & 
    {\bf ho}{\cat M}_2 \ultwocell<\omit>{<0>\eps}
  }$
\end{center}
\item
The {\it total right derived} functor of $u$, denoted $({\bf R}u, \nu)$, is the right Kan extension $({\bf R}_{\gamma_{{\cat M}_1}} (\gamma_{{\cat M}_2} u), \nu)$ of $\gamma_{{\cat M}_2} u$ along $\gamma_{{\cat M}_1}$
\begin{center}
$\xymatrix {
    {\cat M}_1 \ar[r]^u \ar[d]_{\gamma_{_{{\cat M}_1}}} \drtwocell<\omit>{<0>\eta} & 
    {\cat M}_2 \ar[d]^{\gamma_{_{{\cat M}_2}}} \\
    {\bf ho}{\cat M}_1 \ar[r]_{{\bf R}u} & 
    {\bf ho}{\cat M}_2
  }$
\end{center}
\end{enumerate}
\end{defn}

The total left and derived functors $({\bf L}u, \eps), ({\bf R}u, \nu)$ are defined in terms of the localization functors $\gamma_{{\cat M}_1}, \gamma_{{\cat M}_2}$ and therefore will not change if we replace in the definition ${\cat W}_1$, ${\cat W}_2$ with their saturations $\overline{\cat W}_1$, $\overline{\cat W}_2$.

Note that if $u({\cat W}_1) \subset \overline{\cat W}_2$ then ${\bf L}u = {\bf R}u = {\bf ho}u$.

\section{An existence criterion for total derived functors}

\begin{thm}
\label{thm:generalexistenceleftderived}
Given three categories with weak equivalences $({\cat M}^{'}_1$, ${\cat W}^{'}_1)$, $({\cat M}_1$, ${\cat W}_1)$ and $({\cat M}_2$, ${\cat W}_2)$ and two functors $\xymatrix {{\cat M}^{'}_1 \ar[r]^t & {\cat M}_1 \ar[r]^u & {\cat M}_2}$.
\begin{enumerate}
\item
If $t$ is a left approximation and $ut$ preserves weak equivalences, then $u$ admits a total left derived functor $({\bf L}u, \eps)$. The natural map $\eps \colon ({\bf L}u)(tA^{'}) \Ra utA^{'}$ is an isomorphism in objects $A^{'}$ of ${\cat M}^{'}_1$.
\item
If $t$ is a right approximation and $ut$ preserves weak equivalences, then $u$ admits a total right derived functor $({\bf R}u, \eta)$. The natural map $\eta \colon utA^{'} \Ra ({\bf R}u)(tA^{'})$ is an isomorphism in $A^{'}$.
\end{enumerate}
\end{thm}

\begin{proof}
We only prove (1). Pick a regular left cleavage ${\cat C} = (C, D, p, i, j, \sigma)$ along $t$. 
\begin{equation}
\label{eqn:leftcleavage}
\xymatrix { 
A \ar[rr]^f && B \\
tC(A) \ar[u]^{p_A}_\sim \ar[r]_{ti(f)} & 
tD(f) \ar[ur]^{\sigma(f)} &
tC(B) \ar[l]_\sim^{tj(f)} \ar[u]_{p_B}^\sim
}
\end{equation}

Recall that the functor ${\bf s}_{\cat C}$ in the proof of the Approximation \theoremref{thm:approxthm} is a quasi-inverse of ${\bf ho}t$, and define ${\bf L}u = {\bf ho}(ut) \comp {\bf s}_{\cat C}$.

Let us spell out in detail the functor ${\bf L}u$. For an object $A$ of ${\cat M}_1$, we have $({\bf L}u) (A) = utC(A)$. On maps $f \colon A \ra B$ of ${\cat M}_1$, we have ${\bf L}u(f) = (utj(f))^{-1} uti(f)$.

We define $\eps \colon ({\bf L}u)(A) \ra uA$ as $u(p_A)$. The commutativity of the diagram (\ref{eqn:leftcleavage}) implies that $\eps \colon ({\bf L}u) \gamma_{{\cat M}_1} \Ra \gamma_{{\cat M}_2} u$ is a natural map.

We need to show that the pair $({\bf L}u, \eps)$ is {\it terminal} among pairs $(\Lambda, \lambda)$ where $\Lambda \colon {\bf ho} {\cat M}_1 \ra {\bf ho} {\cat M}_2$ is a functor and $\lambda \colon \Lambda \gamma_{{\cat M}_1} \Ra \gamma_{{\cat M}_2} u$ is a natural transformation. For any object $A$ of ${\cat M}_1$ the sequence of full maps in ${\bf ho}{\cat M}_2$ and their inverses
\begin{equation}
\label{eqn:terminalproperty}
\xymatrix { 
    \Lambda(A) \ar@{-->}[rr]^{\lambda(A)} & & uA \\
    \Lambda tC(A)  \ar[u]^-{\Lambda (p_A)}_-\sim \ar[rr]^-{\lambda (tC(A))} & & utC(A) = {\bf L}u(A) \ar@{-->}[u]_{u(p_A) = \eps(A)}
  }
\end{equation}
defines a map $\delta \colon \Lambda(A) \ra {\bf L}u(A)$. For maps $f \colon A \ra B$ of ${\cat M}_1$ we have a commutative diagram in ${\bf ho}{\cat M}_2$
\begin{equation}
\label{eqn:terminalproperty2}
\xymatrix { 
    \Lambda(A) \ar[rr]^{\Lambda(f)} && \Lambda(B) \\
    \Lambda tC(A) \ar[u]^-{\Lambda (p_A)}_-\sim \ar[r]^{\Lambda ti(f)} \ar[d]_-{\lambda (tC(A))} &
    \Lambda tD(f) \ar[ru]^-{\Lambda \sigma(f)} \ar[d]_-{\lambda (tD(f))} &
    \Lambda tC(B) \ar[u]_-{\Lambda (p_B)}^-\sim \ar[l]_{\Lambda tj(f)}^\sim \ar[d]^-{\lambda (tC(B))} \\
    utC(A) \ar[r]_{uti(f)} & utD(f) & utC(B) \ar[l]^{utj(f)}_\sim    
  }
\end{equation}
where $utj(f)$ is a weak equivalence since $j(f)$ is. The commutativity of $(\ref{eqn:terminalproperty2})$ implies that $\delta \colon \Lambda \Ra {\bf L}u$ is natural in maps of ${\cat M}_1$, therefore natural in maps of ${\bf ho}{\cat M}_1$. 

The commutativity of diagram (\ref{eqn:terminalproperty}) shows that we have $\eps \star \delta = \lambda$. 

The natural map $\eps: ({\bf L}u)(tA^{'}) \Ra utA^{'}$ is an isomorphism for {\it regular} objects $A^{'} \eps {\cat M}^{'}_1$, therefore for any object $A^{'}  \eps {\cat M}^{'}_1$. To see that, note that for $A^{'}$ regular the map $\eps: ({\bf L}u)(tA^{'}) \Ra utA^{'}$ can be identified with $1_{utA^{'}}$. Furthermore, our left cleavage was assumed to be regular so any object in ${\cat M}^{'}_1$ is isomorphic in ${\bf ho}{\cat M}^{'}_1$ to a regular object $A^{'}$.
\end{proof}

\begin{cor}
\label{cor:generalexistenceleftderived}
\mbox{}
\begin{enumerate}
\item
If $t, u$ are as in \theoremref{thm:generalexistenceleftderived} (1) and ${\bf s}$ denotes any quasi-inverse of ${\bf ho}t$, then ${\bf ho}(ut) \, {\bf s}$ is naturally isomorphic to ${\bf L}u$.
\item
If $t, u$ are as in \theoremref{thm:generalexistenceleftderived} (2) and ${\bf s}$ denotes any quasi-inverse of ${\bf ho}t$, then ${\bf ho}(ut) \, {\bf s}$ is naturally isomorphic to ${\bf R}u$. $\square$
\end{enumerate}
\end{cor}

\section{The abstract Quillen adjunction property}

We next state an adjunction property of total derived functors. 
\begin{thm}[Abstract Quillen adjunction]
\label{thm:generalexistencetotalderivedadjoint}
Given four categories with weak equivalences $({\cat M}^{'}_1$, ${\cat W}^{'}_1)$, $({\cat M}_1$, ${\cat W}_1)$, $({\cat M}^{'}_2$, ${\cat W}^{'}_2)$, $({\cat M}_2$, ${\cat W}_2)$ and four functors
\begin{center}
$\xymatrix {{\cat M}^{'}_1 \ar[r]^{t_1} & {\cat M}_1 \ar@<2pt>[r]^{u_1} & {\cat M}_2 \ar@<2pt>[l]^{u_2} & {\cat M}^{'}_2 \ar[l]_{t_2} }$
\end{center}
where 
\begin{enumerate}
\item
$u_1 \dashv u_2$ is an adjoint pair
\item
$t_1$ is a left approximation, $t_2$ is a right approximation
\item
$u_1t_1$ and $u_2t_2$ preserve weak equivalences
\end{enumerate}
then ${\bf L}u_1 \dashv {\bf R}u_2$ is a naturally adjoint pair 
\begin{center}
$\xymatrix {{\bf ho}{\cat M}_1 \ar@<2pt>[r]^{{\bf L}u_1} & {\bf ho}{\cat M}_2 \ar@<2pt>[l]^{{\bf R}u_2} }$
\end{center}
If additionally
\begin{enumerate}
\item[(4)]
(resp. (4r), resp. (4l)).
For any objects $A^{'} \eps {\cat M}^{'}_1$, $B^{'} \eps {\cat M}^{'}_2$, a map $u_1t_1A^{'} \ra t_2B^{'}$ is a weak equivalence iff (resp. if, resp. only if) its adjoint $t_1A^{'} \ra u_2t_2B^{'}$ is a weak equivalence
\end{enumerate}
then ${\bf L}u_1$ and ${\bf R}u_2$ are inverse equivalences of categories (resp. ${\bf R}u_2$ is fully faithful, resp ${\bf L}u_1$ is fully faithful).
\end{thm}

This theorem suggests the following
\index{Quillen!abstract adjunction}
\index{Quillen!abstract equivalence}
\begin{defn}
\label{defn:abstractQuillenadjoint}
We will call the functors $u_1, u_2$ satisfying the properties (1), (2), (3) of \theoremref{thm:generalexistencetotalderivedadjoint} an {\it abstract Quillen adjoint} pair with respect to $t_1, t_2$. If the additional property (4) is satisfied, we will call $u_1, u_2$ an {\it abstract Quillen equivalence} pair with respect to $t_1, t_2$.
\end{defn}

There is a very nice, conceptual proof due to Georges Maltsiniotis \cite{Maltsiniotis3} of \theoremref{thm:generalexistencetotalderivedadjoint}, using the universal property in the definition of total derived functors and \theoremref{thm:generalexistenceleftderived}. 

In this text however, we will let \theoremref{thm:generalexistencetotalderivedadjoint} be a consequence of the theorem below. In preparation, let us introduce more notations and definitions. Given a diagram of functors
\begin{equation}
\label{eqn:partialadj}
\xymatrix{
    {\cat A} \ar[r]^{v_1} \ar[d]_{t_1} & {\cat B} \\
    {\cat C} & {\cat D} \ar[l]^{v_2} \ar[u]_{t_2}
  }
\end{equation}
a {\it partial adjunction} \index{functor!partially adjoint} between $v_1, v_2$ with respect to $t_1, t_2$ is a bijection
\begin{center}
$\zeta \colon Hom_{{\cat B}}(v_1A, t_2D) \cong Hom_{{\cat C}}(t_1A, v_2D)$ 
\end{center}
natural in objects $A \eps {\cat A}_1, D \eps {\cat D}$. Whenever we say that $v_1, v_2$ is an adjoint pair with respect to $t_1, t_2$ we refer to a particular bijection $\zeta$. Note that $\zeta^{-1}$ defines a partial adjunction between $t_1, t_2$ with respect to $v_1, v_2$. 

If in addition we have adjoint pairs $t_1 \dashv s_1$, $s_2 \dashv t_2$ 
\begin{equation}
\label{eqn:partialadj2}
\xymatrix{
    {\cat A} \ar[r]^{v_1} \ar@<-2pt>[d]_{t_1} & {\cat B} \ar@<-2pt>[d]_{s_2} \\
    {\cat C} \ar@<-2pt>[u]_{s_1} & {\cat D} \ar[l]^{v_2} \ar@<-2pt>[u]_{t_2}
  }
\end{equation}
then the partial adjunctions of $v_1, v_2$ with respect to $t_1, t_2$ are in a one to one correspondence with adjunctions $s_2v_1 \dashv s_1v_2$. In the diagram (\ref{eqn:partialadj2}), assuming that $s_1, t_1$ and $s_2, t_2$ are inverse equivalences of categories, then the partial adjuctions of $v_1, v_2$ with respect to $t_1, t_2$ are in a one to one correspondence with adjunctions $s_2v_1 \dashv s_1v_2$, and furthermore in a one to one correspondence with adjunctions $v_1s_1 \dashv v_2 s_2$.

In a diagram of the form
\begin{equation}
\label{eqn:partialadj3}
\xymatrix{
    {\cat A}^{'} \ar[r]^{a} & {\cat A} \ar[r]^{v_1} \ar[d]_{t_1} & {\cat B} & \\
    & {\cat C} & {\cat D} \ar[l]^{v_2} \ar[u]_{t_2} & {\cat D}^{'} \ar[l]_{d}
  }
\end{equation}
a partial adjunction between $v_1, v_2$ with respect to $t_1, t_2$ will induce a partial adjunction between $v_1a, v_2d$ with respect to $t_1a, t_2d$, given by 
\begin{center}
$\zeta \colon Hom_{{\cat B}}(v_1aA^{'}, t_2dD^{'}) \cong Hom_{{\cat C}}(t_1aA^{'}, v_2dD^{'})$ 
\end{center}

We can now state
\begin{thm}[Abstract Quillen partial adjunction]
\label{thm:generalexistencetotalderivedadjoint2}
Given four categories with weak equivalences $({\cat M}^{'}_1$, ${\cat W}^{'}_1)$, $({\cat M}_1$, ${\cat W}_1)$, $({\cat M}^{'}_2$, ${\cat W}^{'}_2)$, $({\cat M}_2$, ${\cat W}_2)$ and four functors $t_1, t_2, v_1, v_2$
\begin{center}
$\xymatrix{
    {\cat M}^{'}_1 \ar[r]^{v_1} \ar[d]_{t_1} & {\cat M}_2 \\
    {\cat M}_1 & {\cat M}^{'}_2 \ar[l]_{v_2} \ar[u]_{t_2} &
    }$
\end{center}
such that:
\begin{enumerate}
\item
$v_1, v_2$ are partially adjoint with respect to $t_1, t_2$, meaning that there exists a bijection
\begin{center}
$\zeta \colon Hom_{{\cat M}_2}(v_1A^{'}, t_2B^{'}) \cong Hom_{{\cat M}_1}(t_1A^{'}, v_2B^{'})$
\end{center}
natural in $A^{'} \eps {\cat M}^{'}_1$ and $B^{'} \eps {\cat M}^{'}_2$
\item
$t_1$ is a left approximation, $t_2$ is a right approximation
\item
$v_1$ and $v_2$ preserve weak equivalences
\end{enumerate}

Then ${\bf ho}v_1$, ${\bf ho}v_2$ are naturally partial adjoint with respect to ${\bf ho}t_1$, ${\bf ho}t_2$. Equivalently, denote ${\bf s}_i$ a quasi-inverse of ${\bf ho}t_i$, and let ${\bf V}_i = {\bf ho}(v_i) {\bf s}_i$ for $i = 1, 2$. Then ${\bf V}_1 \dashv {\bf V}_2$
\begin{center}
$\xymatrix {{\bf ho}{\cat M}_1 \ar@<2pt>[r]^{{\bf V}_1} & {\bf ho}{\cat M}_2 \ar@<2pt>[l]^{{\bf V}_2}}$
\end{center}
is a naturally adjoint pair.

If additionally
\begin{enumerate}
\item[(4)]
(resp. (4r), resp. (4l)).
For any objects $A^{'} \eps {\cat M}^{'}_1$, $B^{'} \eps {\cat M}^{'}_2$, a map $v_1A^{'} \ra t_2B^{'}$ is a weak equivalence iff (resp. if, resp. only if) its partial adjoint $t_1A^{'} \ra v_2B^{'}$ is a weak equivalence
\end{enumerate}
then ${\bf V}_1$ and ${\bf V}_2$ are inverse equivalences of categories (resp. ${\bf V}_2$ is fully faithful, resp ${\bf V}_1$ is fully faithful).
\end{thm}

The following definition is suggested:
\index{Quillen!abstract partial adjunction}
\index{Quillen!abstract partial equivalence}
\begin{defn}
\label{defn:abstractQuillenadjoint2}
We will call the functors $v_1, v_2$ satisfying the properties (1), (2), (3) of \theoremref{thm:generalexistencetotalderivedadjoint2} an {\it abstract Quillen partially adjoint} pair with respect to $t_1, t_2$. If the additional property (4) is satisfied, we will call $v_1, v_2$ an {\it abstract Quillen partial equivalence} pair with respect to $t_1, t_2$.
\end{defn}


\begin{proof}[Proof of \theoremref{thm:generalexistencetotalderivedadjoint} assuming \theoremref{thm:generalexistencetotalderivedadjoint2}]
Since $u_1 \dashv u_2$ is an adjoint pair, we see that $v_1 = u_1t_1$, $v_2 = u_2t_2$ is partially adjoint with respect to $t_1, t_2$. From \corollaryref{cor:generalexistenceleftderived}, we have natural isomorphisms ${\bf L}u_1 \cong {\bf ho}(u_1t_1) \, {\bf s}_1 = {\bf V}_1$ and ${\bf R}u_2 \cong {\bf ho}(u_2t_2) \, {\bf s}_2 = {\bf V}_2$. The statement now follows.
\end{proof}

\begin{proof}[Proof of \theoremref{thm:generalexistencetotalderivedadjoint2}]
\mbox{} 
If we can prove the conclusion for a particular choice of ${\bf s}_1 \colon {\bf ho}{\cat M}_1 \ra {\bf ho}{\cat M}^{'}_1$ and ${\bf s}_2 \colon {\bf ho}{\cat M}_2 \ra {\bf ho}{\cat M}^{'}_2$, then the conclusion follows for any ${\bf s}_1$ and ${\bf s}_2$.

We pick a regular left cleavage ${\cat C} = (C, D, p, i, j, \sigma)$ along $t_1$, and a regular right cleavage ${\cat R} = (R, S, i, p^{'}, q^{'}, \tau)$ along $t_2$. As in the proof of the Approximation \theoremref{thm:approxthm}, we get a quasi-inverse ${\bf s}_1 = {\bf s}_{\cat C}$ to ${\bf ho}t_1$ and a quasi-inverse ${\bf s}_2 = {\bf s}_{\cat R}$ to ${\bf ho}t_2$. We will work with these particular choices ${\bf s}_1$ and ${\bf s}_2$.

Denote ${\bf V}_1 = {\bf ho}(v_1) {\bf s}_1$ and ${\bf V}_2 = {\bf ho}(v_2) {\bf s}_2$. We will construct natural maps ${\bf \Phi} \colon 1_{{\bf ho}{\cat M}_1} \Ra {\bf V}_2{\bf V}_1$ and ${\bf \Psi} \colon {\bf V}_1{\bf V}_2 \Ra 1_{{\bf ho}{\cat M}_2}$, and show that they are the unit and counit of an adjunction between ${\bf V}_1$ and ${\bf V}_2$. We start by constructing a natural map
\begin{center}
$\overline{\bf \Phi} \colon {\bf ho}t_1 \Ra {\bf ho}(v_2) \, {\bf s}_2 \, {\bf ho}(v_1)$
\end{center}
by defining $\overline{\bf \Phi} (A^{'}) = \zeta i_{v_1A^{'}}$ for any object $A^{'}$ of ${\bf ho}{\cat M}^{'}_1$, where $i_{v_1A^{'}}$ and $\zeta i_{v_1A^{'}}$ are the maps 
\begin{center}
$\xymatrix{
    v_1A^{'} \ar[r]^-{i_{v_1A^{'}}}_-\sim & t_2R(v_1A^{'}) & t_1A^{'} \ar[rrr]^-{\zeta i_{v_1A^{'}}} &&& v_2R(v_1A^{'})
  }$
\end{center}

Given a map $f \colon A^{'} \ra B^{'}$ in ${\cat M}^{'}_1$ we get a commutative diagram
\[
\xymatrix {
  v_1A^{'} \ar[rr]^{v_1f} \ar[d]_-{i_{v_1A^{'}}}^-\sim \ar[dr]^{\tau(v_1f)} && 
  v_1B^{'} \ar[d]^-{i_{v_1B^{'}}}_-\sim \\
  t_2R(v_1A^{'}) &
  t_2S(v_1f) \ar[l]^{t_2q(v_1f)}_\sim \ar[r]_{t_2p(v_1f)} &
  t_2R(v_1B^{'}) 
}
\]
in ${\cat M}_2$, where $q(v_1f)$ is a weak equivalence. Applying the natural bijection $\zeta$ we get a commutative diagram in ${\cat M}_1$
\[
\xymatrix {
  t_1A^{'} \ar[rr]^{t_1f} \ar[d]_-{\zeta i_{v_1A^{'}}} \ar[dr]^{\zeta \tau(v_1f)} &&
  t_1B^{'} \ar[d]^-{\zeta i_{v_1B^{'}}} \\
  v_2R(v_1A^{'}) &
  v_2S(v_1f) \ar[l]^{v_2q(v_1f)}_\sim \ar[r]_{v_2p(v_1f)} &
  v_2R(v_1B^{'})
}
\]
where $v_2q(v_1f)$ is a weak equivalence since $q(v_1f)$ is. The commutativity of the second diagram shows that $\overline{\bf \Phi}$ is natural in maps of ${\cat M}^{'}_1$, and therefore natural in maps of ${\bf ho}{\cat M}^{'}_1$.

Since ${\bf ho}t_1$ and ${\bf s}_1$ are quasi-inverses of each other, the natural map $\overline{\bf \Phi} \colon {\bf ho}t_1 \Ra {\bf ho}(v_2) \, {\bf s}_2 \, {\bf ho}(v_1)$ yields the desired natural map 
\begin{center}
${\bf \Phi} \colon 1_{{\bf ho}{\cat M}_1} \Ra  {\bf ho}(v_2) \, {\bf s}_2 \, {\bf ho}(v_1) \, {\bf s}_1 = {\bf V}_2{\bf V}_1$
\end{center}

We dually construct a natural map
\begin{center}
$\overline{\bf \Psi} \colon {\bf ho}(v_1) \, {\bf s}_1 \, {\bf ho}(v_2) \Ra {\bf ho}t_2$ 
\end{center}
by defining $\overline{\bf \Psi} (B^{'}) = \zeta^{-1} p_{v_2B^{'}}$ for any object $B^{'}$ of ${\bf ho}{\cat M}^{'}_2$, where the maps $p_{v_2B^{'}}$ and $\zeta^{-1} p_{v_2B^{'}}$ are the maps
\begin{center}
$\xymatrix {
    t_1C(v_2B^{'}) \ar[r]^-{p_{v_2B^{'}}}_-\sim & v_2B^{'} &  & v_1C(v_2B^{'}) \ar[rr]^-{\zeta^{-1} p_{v_2B^{'}}} && t_2B^{'}
  }$
\end{center}
The proof that $\overline{\bf \Psi}$ is a natural map is dual to the proof that $\overline{\bf \Phi}$ is a natural map. Since ${\bf ho}t_2$ and ${\bf s}_2$ are quasi-inverses of each other, the natural map $\overline{\bf \Psi}$ yields the desired natural map 
\begin{center}
${\bf \Psi} \colon {\bf V}_1{\bf V}_2 = {\bf ho}(v_1) \, {\bf s}_1 \, {\bf ho}(v_2) \, {\bf s}_2 \Ra 1_{{\bf ho}{\cat M}_2} $
\end{center}

It remains to show that the natural maps ${\bf \Phi}, {\bf \Psi}$ are the unit resp. the counit of an adjunction between the functors ${\bf V}_1, {\bf V}_2$. In other words, we need to prove that the following composites are identities.
\begin{center}
\begin{equation}
\label{eqn:generaladj11}
\xymatrix {
    {\bf V}_1 \ar@{=>}[r]^-{{\bf V}_1 {\bf \Phi}} & {\bf V}_1{\bf V}_2{\bf V}_1 \ar@{=>}[r]^-{{\bf \Psi} {\bf V}_1} & {\bf V}_1
  }
\end{equation}
\begin{equation}
\label{eqn:generaladj12}
\xymatrix {
    {\bf V}_2 \ar@{=>}[r]^-{{\bf \Phi} {\bf V}_2} & {\bf V}_2{\bf V}_1{\bf V}_2 \ar@{=>}[r]^-{{\bf V}_2 {\bf \Psi}} & {\bf V}_2 
  }
\end{equation}
\end{center}

We only prove that (\ref{eqn:generaladj11}) is an identity, since the proof for (\ref{eqn:generaladj12}) is dual. It suffices to prove that (\ref{eqn:generaladj11}) is an identity on objects of the form $A = t_1A^{'}$, with $A^{'} \eps {\cat M}^{'}_1$ regular with respect to the left cleavage along $t_1$.

If $A^{'}$ is a regular object of ${\cat M}^{'}_1$, we have that $C(t_1A^{'}) = A^{'}$ and $p_{tA^{'}} = 1_{t_1A^{'}} \colon t_1C(t_1A^{'}) \ra t_1A^{'}$.

Denote $h = \zeta i_{v_1A^{'}}$. The left cleavage diagram associated to $h$ yields a commutative diagram in ${\cat M}_1$
\[
\xymatrix { 
  t_1A^{'} \ar[rr]^{h = \zeta i_{v_1A^{'}}} &&
  v_2R(v_1A^{'}) \\
  t_1C(t_1A^{'}) = t_1 A^{'} \ar[u]^-{p_{t_1A^{'}} = 1_{t_1A^{'}}} \ar[r]_-{t_1i(h)} &
  t_1D(h) \ar[ur]^{\sigma(h)} &
  t_1C(v_2R(v_1A^{'})) \ar[l]_-\sim^-{t_1j(h)} \ar[u]_-{p_{v_2R(v_1A^{'})}}^\sim 
}
\]

Applying the natural bijection $\zeta^{-1}$ to the diagram we get a commutative diagram in ${\cat M}_2$
\[
\xymatrix { 
  v_1A^{'} \ar[rr]^{\zeta^{-1}h = i_{v_1A^{'}}}_\sim &&
  t_2R(v_1A^{'}) \\
  v_1A^{'} \ar[u]^{1_{v_1A^{'}}} \ar[r]_-{v_1i(h)} &
  v_1D(h) \ar[ur]^{\zeta^{-1}\sigma(h)} &
  v_1C(v_2R(v_1A^{'})) \ar[l]_-\sim^-{v_1j(h)} \ar[u]_-{\zeta^{-1}p_{v_2R(v_1A^{'})}} 
}
\]
In this diagram, $j(h)$ and therefore $v_1j(h)$ are weak equivalences. 

Since $A = t_1A^{'}$, in ${\bf ho}{\cat M}_1$ we can identify $v_1A^{'}$ with ${\bf V}_1A$ and $v_1C(v_2R(v_1A^{'}))$ with ${\bf V}_1{\bf V}_2{\bf V}_1A$. Under this identification the composition $(v_1j(h))^{-1} \comp v_1i(h)$ becomes ${\bf V}_1{\bf \Phi}(A)$, and the composition $(i_{v_1A^{'}})^{-1}\comp (\zeta^{-1}p_{v_2R(v_1A^{'})})$ becomes ${\bf \Psi}{\bf V}_1(A)$. It follows that the composition (\ref{eqn:generaladj11}) is an identity. Dually, (\ref{eqn:generaladj12}) is an identity and we have proved that ${\bf V}_1 \dashv {\bf V}_2$ are adjoint with adjunction unit ${\bf \Phi}$ and counit ${\bf \Psi}$.

For the second part of the theorem, we assume hypothesis (4l) and we will show that ${\bf \Phi}$ is a natural isomorphisms. From \propositionref{prop:adjfunctors}, this will imply that ${\bf V}_1$ is fully faithful. 

For any object $A^{'}$ of ${\bf ho}{\cat M}^{'}_1$, the map
\begin{center}
$\xymatrix{
    v_1A^{'} \ar[r]^-{i_{v_1A^{'}}}_-\sim & t_2R(v_1A^{'}) 
  }$
\end{center}
and therefore from hypothesis (4) the map
\begin{center}
$\xymatrix{
    t_1A^{'} \ar[rrr]^-{\zeta i_{v_1A^{'}}}_-\sim &&& v_2R(v_1A^{'})
  }$
\end{center}
are weak equivalences. It follows that the natural maps $\overline{\bf \Phi}$ and therefore ${\bf \Phi}$ are isomorphisms. A dual proof shows that hypothesis (4l) implies that ${\bf \Psi}$ is a natural isomorphism, and therefore ${\bf V}_2$ is fully faithful. 

If hypothesis (4) is satisfied, then both ${\bf \Phi}$ and ${\bf \Psi}$ are isomorphisms, therefore ${\bf V}_1, {\bf V}_2$ are inverse equivalences of categories (cf. \propositionref{prop:adjfunctors}).
\end{proof}



